\chapter{Generics}\label{ch:generics}
\index{generics}

Up to this point, every type annotation in our structs, enums, and functions
has been either a concrete type like \lstinline|Int| or the catch-all
\lstinline|any|. Concrete types are precise but inflexible; \lstinline|any|
is flexible but tells the reader nothing about intent. Generics occupy
the middle ground: they let you write code that is \emph{parameterised}
over types, so that the shape of the data is documented even when the
specific type is left open.

In Lattice, generics are a \emph{documentation and design} mechanism.
The language parses generic type parameters on functions, structs, and
enums, and it supports rich type expressions like \lstinline|Map<String, Int>|
and \lstinline|Fn(Int) -> Bool|. At runtime, however, Lattice is
dynamically typed---type parameters are erased and values flow freely.
This chapter shows you how to use generics to write clearer, more
intention-revealing code.

%% ─────────────────────────────────────────────────────────────────────
\section{Generic Functions}\label{sec:generics-functions}
\index{generics!functions}
\index{function!generic}

A generic function declares one or more type parameters in angle brackets
between the function name and the parameter list:

\begin{lstlisting}[language=Lattice, caption={A generic identity function}]
fn identity<T>(x: T) -> T {
    return x
}

print(identity(42))       // 42
print(identity("hello"))  // hello
print(identity(true))     // true
\end{lstlisting}

The \lstinline|<T>| after the function name introduces a type parameter
named \lstinline|T|. We then use \lstinline|T| in the parameter type and
the return type to express that ``whatever type goes in, the same type comes
out.''

\begin{defnbox}{Type Parameter}
A \emph{type parameter} is a placeholder name (conventionally a single
uppercase letter like \lstinline|T|, \lstinline|U|, or \lstinline|K|)
that stands for a type to be specified later. It appears between angle
brackets in a declaration: \lstinline|fn foo<T>(x: T) -> T|.
\end{defnbox}

\subsection{Multiple Type Parameters}

Functions can accept more than one type parameter:

\begin{lstlisting}[language=Lattice, caption={Multiple type parameters}]
fn pair<A, B>(first: A, second: B) -> any {
    return [first, second]
}

let result = pair(1, "two")
print(result)  // [1, two]
\end{lstlisting}

The names \lstinline|A| and \lstinline|B| signal to the reader that the
first and second arguments may be different types.

\subsection{When to Use Generic Functions}

Generic type parameters shine when a function's logic is independent of the
specific type. Classic examples include:

\begin{lstlisting}[language=Lattice, caption={Generic utility functions}]
fn first_of<T>(items: [T]) -> T {
    return items[0]
}

fn swap<A, B>(a: A, b: B) -> any {
    return [b, a]
}

fn repeat<T>(value: T, count: Int) -> [T] {
    flux result = []
    for i in 0..count {
        result.push(value)
    }
    return result
}

print(first_of([10, 20, 30]))    // 10
print(swap("hello", 42))         // [42, hello]
print(repeat("ha", 3))           // [ha, ha, ha]
\end{lstlisting}

Without generics, these functions would use \lstinline|any| for every
parameter---which is \emph{correct} but says nothing about the relationship
between inputs and outputs. The generic signatures communicate intent:
\lstinline|first_of<T>| takes an array of \lstinline|T| and returns a
\lstinline|T|, so you know the return type matches the element type.

\begin{notebox}{Generics Are Erased at Runtime}
Lattice parses generic type parameters and stores them in the AST, but
they are \emph{not enforced} at runtime. The runtime treats all values
as dynamically typed. This means generics serve as \emph{documentation}
and \emph{design intent}---they guide human readers and future tooling,
but the compiler does not yet reject type mismatches. Think of them as
structured comments with syntactic validation.
\end{notebox}

\subsection{Generic Functions with Constraints}

Although Lattice does not enforce type constraints today, you can document
expected capabilities in comments:

\begin{lstlisting}[language=Lattice, caption={Documenting type expectations}]
// T must support the `+` operator
fn sum_all<T>(items: [T], zero: T) -> T {
    flux total = zero
    for item in items {
        total = total + item
    }
    return total
}

print(sum_all([1, 2, 3, 4], 0))                // 10
print(sum_all([1.5, 2.5, 3.0], 0.0))           // 7.0
print(sum_all(["a", "b", "c"], ""))             // abc
\end{lstlisting}

The function works for any type that supports addition. The generic
parameter \lstinline|T| communicates that the accumulator, the elements,
and the result are all the same type.

%% ─────────────────────────────────────────────────────────────────────
\section{Generic Structs}\label{sec:generics-structs}
\index{generics!structs}
\index{struct!generic}

Structs can also declare type parameters:

\begin{lstlisting}[language=Lattice, caption={A generic Box struct}]
struct Box<T> {
    value: T
}

let int_box = Box { value: 42 }
let str_box = Box { value: "hello" }

print(int_box.value)  // 42
print(str_box.value)  // hello
\end{lstlisting}

The \lstinline|<T>| after the struct name introduces a type parameter.
The field \lstinline|value| is annotated with \lstinline|T|, signalling
that a \lstinline|Box| holds exactly one value of a uniform type.

\begin{tipbox}{Instantiation Does Not Require Type Arguments}
When you write \lstinline|Box { value: 42 }|, you do not need to specify
\lstinline|Box<Int> { value: 42 }|. The type parameter exists in the
\emph{definition} to document the struct's generic nature; at construction
time, the value itself determines the concrete type.
\end{tipbox}

\subsection{Multi-Parameter Structs}

\begin{lstlisting}[language=Lattice, caption={A struct with two type parameters}]
struct Pair<A, B> {
    first: A,
    second: B
}

let entry = Pair { first: "name", second: 42 }
print(entry.first)   // name
print(entry.second)  // 42
\end{lstlisting}

The \lstinline|Pair<A, B>| struct makes it clear that the two fields may
hold different types. Contrast this with:

\begin{lstlisting}[language=Lattice]
struct Pair { first: any, second: any }
\end{lstlisting}

Both definitions compile identically, but the generic version tells the
reader ``these are distinct type roles, not arbitrary values.''

\subsection{Realistic Generic Structs}

Let's build a generic key--value entry and a lightweight container:

\begin{lstlisting}[language=Lattice, caption={Generic key-value entry}]
struct Entry<K, V> {
    key: K,
    value: V
}

fn make_entry<K, V>(key: K, value: V) -> Entry {
    return Entry { key: key, value: value }
}

let user_entry = make_entry("alice", 30)
print(user_entry.key)    // alice
print(user_entry.value)  // 30
\end{lstlisting}

\begin{lstlisting}[language=Lattice, caption={A generic wrapper with a label}]
struct Tagged<T> {
    tag: String,
    data: T
}

let measurement = Tagged { tag: "temperature", data: 23.5 }
let flag = Tagged { tag: "active", data: true }

print(measurement.tag + ": " + to_string(measurement.data))  // temperature: 23.5
print(flag.tag + ": " + to_string(flag.data))                // active: true
\end{lstlisting}

\subsection{Generic Structs with Traits}

Generic structs work seamlessly with trait implementations:

\begin{lstlisting}[language=Lattice, caption={Traits on generic structs}]
struct Box<T> {
    value: T
}

trait Describable {
    fn describe(self: any) -> String
}

impl Describable for Box {
    fn describe(self: any) -> String {
        return "Box(" + to_string(self.value) + ")"
    }
}

let b = Box { value: 42 }
print(b.describe())  // Box(42)

let s = Box { value: "world" }
print(s.describe())  // Box(world)
\end{lstlisting}

The impl block says \lstinline|impl Describable for Box| (without type
parameters)---because at runtime, all \lstinline|Box| instances share the
same type name regardless of what \lstinline|T| was.

%% ─────────────────────────────────────────────────────────────────────
\section{Generic Enums and Trait Implementations}\label{sec:generics-enums}
\index{generics!enums}
\index{enum!generic}

Enums are perhaps the most natural fit for generics. The
\lstinline|Option<T>| and \lstinline|Result<T, E>| patterns from
\cref{ch:enums} become far more expressive with type parameters:

\begin{lstlisting}[language=Lattice, caption={A generic Option enum}]
enum Option<T> {
    Some(T),
    None
}

let found = Option::Some(42)
let missing = Option::None

print(found)    // Option::Some(42)
print(missing)  // Option::None
\end{lstlisting}

The type parameter \lstinline|<T>| makes the intent unmistakable:
\lstinline|Some| carries a value of type \lstinline|T|, and \lstinline|None|
carries nothing.

\subsection{Generic Result}

\begin{lstlisting}[language=Lattice, caption={A generic Result enum}]
enum Result<T, E> {
    Ok(T),
    Err(E)
}

fn divide(a: Int, b: Int) -> Result {
    if b == 0 {
        return Result::Err("division by zero")
    }
    return Result::Ok(a / b)
}

let r = divide(10, 3)
match r {
    Result::Ok(v)  => print("result: " + to_string(v)),
    Result::Err(e) => print("error: " + e),
    _ => print("?")
}
// Output: result: 3
\end{lstlisting}

With \lstinline|Result<T, E>|, the definition documents that the success
payload is one type and the error payload is another. The two-parameter
signature mirrors the Rust and Haskell \texttt{Either} type.

\subsection{Recursive Generic Enums}

The \lstinline|Tree| example from \cref{sec:enum-recursive} becomes more
self-documenting with a type parameter:

\begin{lstlisting}[language=Lattice, caption={A generic binary tree}]
enum Tree<T> {
    Leaf(T),
    Branch(any, any)
}

fn tree_sum(t: any) -> Int {
    match t {
        Tree::Leaf(v)      => v,
        Tree::Branch(l, r) => tree_sum(l) + tree_sum(r),
        _ => 0
    }
}

let tree = Tree::Branch(
    Tree::Branch(Tree::Leaf(1), Tree::Leaf(2)),
    Tree::Leaf(3)
)
print(tree_sum(tree))  // 6
\end{lstlisting}

\begin{notebox}{Branch Payloads and Self-Reference}
You might notice that \lstinline|Branch| uses \lstinline|any| for its
children rather than \lstinline|Tree<T>|. This is because Lattice's type
system does not yet resolve recursive type references within enum payloads.
Using \lstinline|any| for the recursive positions is the idiomatic approach
until the type system evolves.
\end{notebox}

%% ─────────────────────────────────────────────────────────────────────
\section{Type Expressions}\label{sec:generics-type-expressions}
\index{type expression}
\index{type annotation}

Beyond simple names like \lstinline|Int| and \lstinline|String|, Lattice
supports a rich syntax for \emph{type expressions}---the annotations you
write after colons in parameter lists and struct fields.

\subsection{Named Types}

The most basic type expression is a single name:

\begin{lstlisting}[language=Lattice]
fn greet(name: String) -> String {
    return "Hello, " + name
}
\end{lstlisting}

Built-in named types include \lstinline|Int|, \lstinline|Float|,
\lstinline|Bool|, \lstinline|String|, \lstinline|Array|, \lstinline|Map|,
\lstinline|Set|, \lstinline|Tuple|, \lstinline|Channel|, \lstinline|Buffer|,
\lstinline|Iterator|, \lstinline|Ref|, and \lstinline|Fn|. Custom struct
and enum names are also valid.

\subsection{Parameterised Named Types}

Named types can take type arguments in angle brackets:

\begin{lstlisting}[language=Lattice, caption={Parameterised type expressions}]
struct Config {
    headers: Map<String, String>,
    ports: Array<Int>,
    handler: Fn
}
\end{lstlisting}

\lstinline|Map<String, String>| says ``a map from strings to strings.''
\lstinline|Array<Int>| says ``an array of integers.'' The parser recognises
the angle brackets and stores the type arguments in the AST (see
\filename{src/parser.c}, the \texttt{parse\_type\_expr} function).

\subsection{Array Types}

Lattice supports a bracket syntax for array types:

\begin{lstlisting}[language=Lattice, caption={Array type annotation}]
fn sum(nums: [Int]) -> Int {
    flux total = 0
    for n in nums {
        total = total + n
    }
    return total
}

print(sum([1, 2, 3, 4, 5]))  // 15
\end{lstlisting}

The annotation \lstinline|[Int]| is shorthand for ``an array whose elements
are integers.'' It is syntactically distinct from \lstinline|Array<Int>| but
carries the same meaning.

\subsection{Function Types}
\index{function type}

Function types describe the signature of a callable value:

\begin{lstlisting}[language=Lattice, caption={Function type annotations}]
fn apply(f: Fn(Int) -> Int, x: Int) -> Int {
    return f(x)
}

let double = |n| { return n * 2 }
print(apply(double, 5))  // 10
\end{lstlisting}

The type \lstinline|Fn(Int) -> Int| means ``a function that takes an
\lstinline|Int| and returns an \lstinline|Int|.'' Multi-parameter function
types list all parameter types:

\begin{lstlisting}[language=Lattice, caption={Multi-parameter function types}]
fn apply_binary(f: Fn(Int, Int) -> Int, a: Int, b: Int) -> Int {
    return f(a, b)
}

let add = |x, y| { return x + y }
print(apply_binary(add, 3, 4))  // 7
\end{lstlisting}

\subsection{Combining Type Expressions}

These forms compose naturally. Here is a function that takes a
transformation function and an array, and returns a new array:

\begin{lstlisting}[language=Lattice, caption={Combining function and array types}]
fn transform<T, U>(items: [T], f: Fn(T) -> U) -> [U] {
    flux result = []
    for item in items {
        result.push(f(item))
    }
    return result
}

let names = ["alice", "bob", "charlie"]
let lengths = transform(names, |s| { return len(s) })
print(lengths)  // [5, 3, 7]
\end{lstlisting}

The signature \lstinline|fn transform<T, U>(items: [T], f: Fn(T) -> U) -> [U]|
is dense but readable: it takes an array of \lstinline|T|, a function from
\lstinline|T| to \lstinline|U|, and returns an array of \lstinline|U|.
This is the classic ``map'' operation expressed with full type information.

\begin{tipbox}{Type Expressions as Documentation}
Even though Lattice does not enforce type expressions at runtime, they serve
as \emph{machine-readable documentation}. Future tooling---such as the
language server, linter, and doc generator---can leverage these annotations
to provide better auto-complete, hover information, and static analysis.
Write type expressions for your future self and your teammates.
\end{tipbox}

\subsection{Phase Annotations in Types}

Type expressions can also include phase qualifiers from the phase system:

\begin{lstlisting}[language=Lattice, caption={Phase-annotated types}]
fn freeze_point(p: ~Point) -> *Point {
    return freeze p
}
\end{lstlisting}

The tilde (\lstinline|~|) denotes a fluid (mutable) value; the asterisk
(\lstinline|*|) denotes a crystal (frozen) value. These can be combined
with any type expression, including generic ones. Phase annotations are
covered in depth in \cref{ch:phases-explained}.

%% ─────────────────────────────────────────────────────────────────────
\section{Practical Patterns with Generics}\label{sec:generics-patterns}
\index{generics!patterns}

Let's explore several design patterns that benefit from generic annotations.

\subsection{The Wrapper Pattern}

A common need is to wrap a value with metadata:

\begin{lstlisting}[language=Lattice, caption={A timestamped wrapper}]
struct Timestamped<T> {
    value: T,
    created_at: Int
}

fn stamp<T>(value: T, time: Int) -> Timestamped {
    return Timestamped { value: value, created_at: time }
}

let record = stamp("measurement: 23.5C", 1700000000)
print(record.value)       // measurement: 23.5C
print(record.created_at)  // 1700000000
\end{lstlisting}

The type parameter \lstinline|T| signals that \lstinline|Timestamped| is
agnostic about what it wraps. It could hold a string, an integer, or another
struct.

\subsection{The Builder Pattern}

Generic structs pair well with method chaining:

\begin{lstlisting}[language=Lattice, caption={A generic builder}]
struct Builder<T> {
    parts: [T],
    count: Int
}

trait Buildable {
    fn add(self: any, item: any) -> any
    fn items(self: any) -> any
}

impl Buildable for Builder {
    fn add(self: any, item: any) -> any {
        flux new_parts = []
        for i in 0..self.count {
            new_parts.push(self.parts[i])
        }
        new_parts.push(item)
        return Builder { parts: new_parts, count: self.count + 1 }
    }

    fn items(self: any) -> any {
        return self.parts
    }
}

let list = Builder { parts: [], count: 0 }
let result = list.add("alpha").add("beta").add("gamma")
print(result.items())  // [alpha, beta, gamma]
\end{lstlisting}

\subsection{Higher-Order Generic Functions}

Generic annotations make higher-order functions self-documenting:

\begin{lstlisting}[language=Lattice, caption={Generic filter and reduce}]
fn filter<T>(items: [T], predicate: Fn(T) -> Bool) -> [T] {
    flux result = []
    for item in items {
        if predicate(item) {
            result.push(item)
        }
    }
    return result
}

fn reduce<T, U>(items: [T], initial: U, combine: Fn(U, T) -> U) -> U {
    flux acc = initial
    for item in items {
        acc = combine(acc, item)
    }
    return acc
}

let nums = [1, 2, 3, 4, 5, 6, 7, 8, 9, 10]

let evens = filter(nums, |n| { return n % 2 == 0 })
print(evens)  // [2, 4, 6, 8, 10]

let total = reduce(nums, 0, |acc, n| { return acc + n })
print(total)  // 55
\end{lstlisting}

The signature of \lstinline|reduce<T, U>| makes the roles crystal clear:
\lstinline|T| is the element type, \lstinline|U| is the accumulator type,
and the \lstinline|combine| function bridges the two.

\subsection{Generic Option Utilities}

Building on the \lstinline|Option| enum from \cref{sec:enum-option-result}:

\begin{lstlisting}[language=Lattice, caption={Generic Option helpers}]
enum Option<T> {
    Some(T),
    None
}

fn map_option<T, U>(opt: Option, f: Fn(T) -> U) -> Option {
    match opt {
        Option::Some(v) => Option::Some(f(v)),
        Option::None    => Option::None,
        _ => Option::None
    }
}

fn unwrap_or<T>(opt: Option, default: T) -> T {
    match opt {
        Option::Some(v) => v,
        Option::None    => default,
        _ => default
    }
}

let name = Option::Some("Alice")
let upper = map_option(name, |s| { return s.to_upper() })
print(unwrap_or(upper, "unknown"))  // ALICE

let missing = Option::None
print(unwrap_or(missing, "anonymous"))  // anonymous
\end{lstlisting}

These utility functions work with any \lstinline|Option|, regardless of
the payload type, because the logic depends only on the variant, not on
\lstinline|T|.

\subsection{Combining Generics, Enums, and Traits}

Here is a pattern that brings all three features together---a generic
``result handler'' that describes and processes outcomes:

\begin{lstlisting}[language=Lattice, caption={Combining generics with enums and traits}]
enum Outcome<T> {
    Success(T),
    Failure(any)
}

struct Report<T> {
    outcome: Outcome,
    label: String
}

trait Summarizable {
    fn summary(self: any) -> String
}

impl Summarizable for Report {
    fn summary(self: any) -> String {
        let status = match self.outcome {
            Outcome::Success(v) => "OK: " + to_string(v),
            Outcome::Failure(e) => "FAIL: " + to_string(e),
            _ => "UNKNOWN"
        }
        return "[" + self.label + "] " + status
    }
}

let good = Report {
    outcome: Outcome::Success(200),
    label: "health check"
}

let bad = Report {
    outcome: Outcome::Failure("timeout"),
    label: "database ping"
}

print(good.summary())  // [health check] OK: 200
print(bad.summary())   // [database ping] FAIL: timeout
\end{lstlisting}

The \lstinline|Report<T>| struct holds an \lstinline|Outcome<T>| enum, and
the \lstinline|Summarizable| trait provides a polymorphic \lstinline|summary|
method. Each piece---generics, enums, traits---contributes a different kind
of abstraction, and they compose cleanly.

%% ─────────────────────────────────────────────────────────────────────
\section{The Future of Generics in Lattice}\label{sec:generics-future}
\index{generics!roadmap}

Lattice's generic system today is best described as \emph{syntax-first}:
the parser fully understands type parameters and type arguments, and the
AST preserves them, but the compiler and runtime do not yet enforce them.
This is a deliberate staged approach.

What works today:
\begin{itemize}
  \item Generic type parameter syntax on functions, structs, and enums.
  \item Type arguments in type expressions (\lstinline|Map<String, Int>|).
  \item Function type expressions (\lstinline|Fn(Int) -> Bool|).
  \item Array type shorthand (\lstinline|[Int]|).
  \item Phase-qualified type expressions.
\end{itemize}

What is planned for future versions:
\begin{itemize}
  \item Compile-time type checking against declared parameters.
  \item Trait bounds on type parameters (\lstinline|<T: Describable>|).
  \item Better error messages when types do not match at call sites.
  \item Tooling support: the language server already reads type annotations
    for hover information (see \filename{src/lsp\_server.c}).
\end{itemize}

\begin{warnbox}{Generics Do Not Prevent Type Errors (Yet)}
Because type parameters are erased at runtime, passing a
\lstinline|String| where an \lstinline|Int| is expected will not produce
a compile-time error---you will get a runtime error when the value is used
incorrectly. Use generics for documentation today, and look forward to
enforcement in future Lattice versions.
\end{warnbox}

Write your generic annotations now. They cost nothing at runtime, they make
your code dramatically easier to understand, and when the type checker
arrives, your code will already be ready.

%% ─────────────────────────────────────────────────────────────────────
\section{Exercises}\label{sec:generics-exercises}

\begin{enumerate}
  \item \textbf{Generic stack.}
    Define a \lstinline|struct Stack<T>| with fields \lstinline|items: [T]|
    and \lstinline|size: Int|. Implement a \lstinline|Pushable| trait with
    \lstinline|push|, \lstinline|pop|, and \lstinline|peek| methods. Test
    it with both integers and strings.

  \item \textbf{Option pipeline.}
    Write three functions: \lstinline|map_option<T, U>|,
    \lstinline|flat_map_option<T, U>| (where \lstinline|f| returns an
    \lstinline|Option|), and \lstinline|or_else<T>| (provides a fallback
    option). Chain them to transform \lstinline|Option::Some(5)| into
    \lstinline|Option::Some("10")| (double, then convert to string).

  \item \textbf{Generic key-value store.}
    Define \lstinline|struct KVStore<K, V>| that wraps a \lstinline|Map|.
    Implement \lstinline|set|, \lstinline|get|, and \lstinline|has| methods.
    Use it to build a phone book mapping names to numbers.

  \item \textbf{Type expression exploration.}
    Write a function whose signature uses all three type expression forms:
    a named type with arguments (\lstinline|Map<String, Int>|), an array
    type (\lstinline|[String]|), and a function type
    (\lstinline|Fn(String) -> Int|). The function should accept a list of
    strings, a scoring function, and return a map from strings to their
    scores.

  \item \textbf{Generic Result chain.}
    Define \lstinline|enum Result<T, E>| and write a function
    \lstinline|and_then<T, U, E>(r: Result, f: Fn(T) -> Result) -> Result|
    that applies \lstinline|f| only if \lstinline|r| is \lstinline|Ok|.
    Chain three \lstinline|and_then| calls to parse a string into an
    integer, double it, and validate that it is positive.
\end{enumerate}

%% ─────────────────────────────────────────────────────────────────────
\section*{What's Next}

With structs, enums, traits, and generics, you now have Lattice's full
toolkit for structuring data and building abstractions. But programs do
not just model data---they \emph{do} things, often many things at once.
In the next part we enter the world of \emph{concurrency}: how Lattice's
structured concurrency model with \lstinline|scope| and \lstinline|spawn|
lets you run tasks in parallel while keeping your code safe and your sanity
intact.
